% literatureReview.tex — Related Literature
% Theoretical foundations for the fee concentration index and exit mechanism

\section{Related Literature}\label{sec:related-literature}

Six papers independently validate the components of $\AT$ and the exit
mechanism that motivates insurance demand.

\cite{capponi2024optimal} model the LP's exit as an optimal stopping problem
with thresholds $p_U^*$ and $p_L^*$. The key comparative static
$\partial p_U^* / \partial \phi > 0$ implies that a higher effective fee rate
raises the exit threshold (LP stays longer). When $\AT$ rises, passive LP's
effective fee rate $\hat{\phi}_i = x_i \cdot \phi_{\text{pool}}$ drops because
$x_i < 1/\PosCount$, lowering $p_U^*$ and triggering earlier exit. This
predicts $\beta_1 > 0$ for sufficiently high concentration.

\cite{capponi2024jit} show that JIT liquidity creates fee share asymmetry: the
passive LP's share is $d_P/(d_P + d_J)$ where $d_J$ is JIT depth---exactly
the $x_k$ component of $\AT$. JIT entry mechanically raises $\AT$ by
concentrating fees in short-lived positions. Our JIT control analysis (32 JIT
positions, 5.3\% share) confirms JIT count has an independent positive effect
on exits ($\beta_{\text{JIT}} = 0.177$, $p = 0.031$), but does not explain
the $\AT$ effect, confirming that $\AT$ captures concentration beyond JIT.

\cite{macrapis2024permissionless} derive the symmetric equilibrium with
$\PosCount$ LPs in free entry, where each earns
$\Pi_i = (1/\PosCount^2) f B_D$ and the equal-share outcome
$x_k = 1/\PosCount$ is the null hypothesis of $\AT$. Our deviation
$\Delta_t = A_T^{\text{real}} - A_T^{1/N}$ tests departure from this
equilibrium directly. The deviation---not $\AT$ in levels---is the correct
treatment variable, since levels confound pool health with concentration
pressure.

\cite{aquilina2024dealers} provide empirical validation: sophisticated LPs
(7\% of addresses) provide 65--85\% of liquidity, use tick ranges
$3.4\times$ tighter, hold positions $7.5\times$ shorter, and capture
75--90\% of fees. This confirms that $x_k$ in $\AT$ is highly skewed. Their
profitability decomposition
$R_{\text{net}} = \text{FeeYield} + \text{IL} - \text{Gas}$ shows that
passive LPs earn negative returns when sophisticated LPs dominate---precisely
when $\Delta_t$ is large.

\cite{bichuch2024flair} formulate a liquidity provider competitiveness metric
(FLAIR) that quantifies the fee rate heterogeneity across positions in a
pool. Their metric captures the same distributional asymmetry that $\AT$
measures---the gap between what passive and sophisticated LPs earn per unit
of liquidity. FLAIR independently confirms that fee share inequality is
persistent and economically significant across Uniswap V3 pools.

\cite{bichuch2024flair} also show that the LP value function is an optimal
stopping problem with implied fee rate
$F(p_x, p_y) = p_y \cdot \ell(p_x/p_y)$ where $\ell$ is instantaneous LVR.
Their Theorem~4.2 establishes that a risk-neutral LP is indifferent across
all stopping times, meaning any deterministic exit driver must come from fee
rate heterogeneity across LPs---exactly what $\AT$ measures. Their
fixed-for-floating fee swap (Corollary~5.5) is structurally identical to a
fee concentration derivative: the LP exchanges volatile realized fees for a
fixed premium.
